% Canonical Paper II source.
The Poisson theorem has two complementary refinements.  Fourier analysis
identifies its boundary energy operator exactly, while the assignment relation
\eqref{eq:hyperbolic-assignment-w} produces explicit arithmetic families.  We
keep their domains separate: a plane wave is a useful bounded Fourier mode but
does not extend continuously to the single ideal point $\infty$, whereas the
$C_0$ theorem of the preceding section does.

\subsection{Fourier representation and the conormal derivative}

Use the Fourier convention
\begin{equation}\label{eq:fourier-convention-x}
 \widehat\varphi(\kappa)
 =\int_{\R}e^{-i\kappa x}\varphi(x)\,dx,
 \qquad
 \varphi(x)=\frac1{2\pi}\int_{\R}
 e^{i\kappa x}\widehat\varphi(\kappa)\,d\kappa.
\end{equation}
Write $|D_x|$ for the multiplier with symbol $|\kappa|$.

\begin{proposition}[Fourier formula for the Poisson extension]
\label{prop:poisson-fourier-formula}
If $\varphi\in\mathcal S(\R)$ and
$U=\mathcal P_{\mu,\lambda}\varphi$, then
\begin{equation}\label{eq:poisson-fourier-x}
 U(x,y)=\frac1{2\pi}\int_{\R}
 \widehat\varphi(\kappa)
 e^{i\kappa x-(\mu/\lambda)|\kappa|y}\,d\kappa.
\end{equation}
Consequently,
\begin{equation}\label{eq:raw-normal-fourier}
 -\partial_yU(\,\cdot\,,0)
 =\frac{\mu}{\lambda}|D_x|\varphi
\end{equation}
in $\mathcal S(\R)$ and in $L^2(\R)$.
\end{proposition}

\begin{proof}
The Fourier transform in $x$ of
$\pi^{-1}h/(x^2+h^2)$ is $e^{-h|\kappa|}$.  In
\eqref{eq:aeg-poisson-kernel-x}, $h=(\mu/\lambda)y$.  The convolution theorem
therefore gives \eqref{eq:poisson-fourier-x}.  Since the data are Schwartz,
one may differentiate under the integral for $y\geq0$ in $L^2$; evaluation
at $y=0$ yields \eqref{eq:raw-normal-fourier}.
\end{proof}

The raw derivative in \eqref{eq:raw-normal-fourier} is not yet the boundary
operator paired with the Riemannian energy.  On the truncated domain
$\{y\geq\varepsilon\}$, the outward $g$-unit normal and induced measure along
$y=\varepsilon$ are
\begin{equation}\label{eq:hyperbolic-boundary-normal-measure}
 \nu=-\lambda y\partial_y,
 \qquad
 d\sigma_g=\frac{dx}{\mu y}.
\end{equation}
Thus
\[
 (\partial_\nu U)d\sigma_g
 =-\frac{\lambda}{\mu}U_y\,dx.
\]
This fixes the conformal conormal scale.

\begin{definition}[Conformal Dirichlet-to-Neumann operator]
\label{def:conformal-DtN}
For $\varphi\in\mathcal S(\R)$, define
\begin{equation}\label{eq:conformal-DtN-x}
 \mathcal N_x\varphi
 :=-\frac{\lambda}{\mu}
   \partial_y(\mathcal P_{\mu,\lambda}\varphi)(\,\cdot\,,0).
\end{equation}
\end{definition}

\begin{theorem}[Dirichlet-to-Neumann and energy identity]
\label{thm:DtN-energy-identity}
For every $\varphi\in\mathcal S(\R)$,
\begin{align}
 \mathcal N_x\varphi&=|D_x|\varphi,
 \label{eq:DtN-is-absolute-D}\\
 \mathcal E_g(\mathcal P_{\mu,\lambda}\varphi)
 &=\int_{\R}\overline{\varphi(x)}\,
       \mathcal N_x\varphi(x)\,dx
  =\frac1{2\pi}\int_{\R}
       |\kappa|\,|\widehat\varphi(\kappa)|^2\,d\kappa.
 \label{eq:poisson-energy-fourier}
\end{align}
For real-valued data the complex conjugate may be omitted.  In the normalized
$X$-coordinate, the same statement reads
\[
 \mathcal N_X=-\partial_y\big|_{y=0}=|D_X|.
\]
\end{theorem}

\begin{proof}
Equation~\eqref{eq:DtN-is-absolute-D} follows immediately from
\eqref{eq:raw-normal-fourier} and
\eqref{eq:conformal-DtN-x}.  To calculate the energy, insert
\eqref{eq:poisson-fourier-x} into
\eqref{eq:hyperbolic-dirichlet-energy} and use Plancherel.  With
$c=\mu/\lambda$,
\begin{align*}
 \mathcal E_g(U)
 &=\frac1{2\pi}\int_{\R}\int_0^\infty
   \left(c\kappa^2+c^{-1}c^2\kappa^2\right)
   |\widehat\varphi(\kappa)|^2e^{-2c|\kappa|y}
   \,dy\,d\kappa\\
 &=\frac1{2\pi}\int_{\R}|\kappa|
   |\widehat\varphi(\kappa)|^2\,d\kappa.
\end{align*}
Plancherel once more identifies this expression with
$\int\overline\varphi\,|D_x|\varphi\,dx$.  Finally, if
$\psi(X)=\varphi((\mu/\lambda)X)$, the $X$-Fourier solution decays as
$e^{-|\xi|y}$, so its raw outward conformal derivative is $|D_X|$.
\end{proof}

The right side of \eqref{eq:poisson-energy-fourier} is the homogeneous
$H^{1/2}(\R)$ quadratic form.  The identity therefore extends by completion
from Schwartz data to boundary data with finite homogeneous
$H^{1/2}$ seminorm, modulo the usual treatment of constants.

\begin{theorem}[Energy-minimizing extension]
\label{thm:poisson-energy-minimizer}
Let $\varphi\in\mathcal S(\R)$ and
$U=\mathcal P_{\mu,\lambda}\varphi$.  For every
$h\in H_{0,\mathrm{Euc}}^1(\HH)$, where this denotes the Euclidean half-plane
Sobolev space,
the function $U+h$ has the same trace as $U$ and
\begin{equation}\label{eq:energy-orthogonal-decomposition}
 \mathcal E_g(U+h)=\mathcal E_g(U)+\mathcal E_g(h).
\end{equation}
Consequently, $U$ is the unique energy minimizer in the affine class
\[
 U+H_{0,\mathrm{Euc}}^1(\HH).
\]
It is also the unique weakly $\Delta_g$-harmonic member of that class.
\end{theorem}

\begin{proof}
The anisotropic energy norm in
\eqref{eq:hyperbolic-dirichlet-energy} is equivalent to the Euclidean
gradient norm because $\mu,\lambda>0$.  For
$h\in C_c^\infty(\HH)$, integration by parts and $\Delta_gU=0$ give
\[
 \int_{\HH}\left(
  \frac{\mu}{\lambda}U_x\overline{h_x}
 +\frac{\lambda}{\mu}U_y\overline{h_y}\right)dx\,dy=0.
\]
Density of $C_c^\infty(\HH)$ in $H_{0,\mathrm{Euc}}^1(\HH)$ extends this orthogonality to
all admissible $h$.  Expanding the squared gradient proves
\eqref{eq:energy-orthogonal-decomposition}.  Equality of energies forces
$\nabla h=0$ almost everywhere; since
$h\in H_{0,\mathrm{Euc}}^1(\HH)\subset L^2(\HH)$,
the constant is zero.  This proves uniqueness of the minimizer.

If $U+h$ is also weakly harmonic, subtract the weak equations for $U+h$ and
$U$, and test the resulting equation for $h$ by $h$ itself, using density.
Then $\mathcal E_g(h)=0$, so $h=0$.
\end{proof}

\subsection{Holomorphic and Fourier families in assignment coordinates}

Equation~\eqref{eq:hyperbolic-dbar-coordinate} immediately gives a broad
family, with an explicit arithmetic presentation.

\begin{proposition}[Assignment presentation of holomorphic fields]
\label{prop:assignment-holomorphic-family}
Let $\Omega\subset\C$ be open and let $\Phi$ be holomorphic on $\Omega$.
On every component of
\[
 \left\{(a,y)\in\R\times(0,\infty):
 y\left(i-\frac{\lambda}{\mu}a\right)\in\Omega\right\},
\]
the field
\begin{equation}\label{eq:assignment-holomorphic-composition}
 F(a,y)=\Phi\!\left(y\left(i-\frac{\lambda}{\mu}a\right)\right)
\end{equation}
is arithmetic holomorphic.  Its real and imaginary parts are
$\Delta_g$-harmonic.
\end{proposition}

\begin{proof}
The argument of $\Phi$ is precisely $w$ by
\eqref{eq:hyperbolic-assignment-w}.  Hence
$\partial_{\bar w}\Phi(w)=0$, and
\eqref{eq:hyperbolic-dbar-coordinate} gives $\dbarAEG F=0$.  Surface
holomorphicity implies Laplace--Beltrami harmonicity by
Section~\ref{sec:surface-operators}.
\end{proof}

Examples include polynomials in $w$ and rational functions $R(w)$ on domains
avoiding their poles.  If all poles of $R$ lie in
$\C\setminus\overline\HH$, then $R(w)$ is globally holomorphic on $\HH$.
For example, for every $\tau>0$,
\begin{equation}\label{eq:explicit-rational-assignment-family}
 R_\tau(w)=\frac1{w+i\tau}
 =\frac1{-(\lambda/\mu)ay+i(y+\tau)}
\end{equation}
is a global arithmetic holomorphic field; its pole $-i\tau$ lies in the lower
half-plane.  Its real and imaginary parts are explicit rational harmonic
functions of $(a,y)$.
The positive-frequency modes are, for $\kappa>0$,
\begin{equation}\label{eq:positive-holomorphic-mode}
 H_\kappa(x,y)
 =e^{i\kappa x-(\mu/\lambda)\kappa y}
 =e^{i(\mu/\lambda)\kappa w}
 =e^{-i\kappa ay-(\mu/\lambda)\kappa y}.
\end{equation}
They are holomorphic.  The decaying modes with negative frequency are
anti-holomorphic.  Formula~\eqref{eq:poisson-fourier-x} is their
superposition.  Individual plane waves have bounded boundary values but no
limit at the compactification point $\infty$, so they are not examples in the
$C_0$ theorem.

\subsection{Assignment-only harmonic fields}

The assignment itself is not harmonic, but one nonlinear rectification of it
is.  In fact there are no others beyond affine recombination.

\begin{theorem}[Assignment-only harmonic classification]
\label{thm:a-only-harmonic-classification}
Let $I\subset\R$ be an open interval, let $h\in C^2(I)$, and put
$\Omega_I=a^{-1}(I)$.  Then
\begin{equation}\label{eq:a-only-laplacian}
 \Delta_g(h\circ a)
 =(\mu^2+\lambda^2a^2)h''(a)
  +2\lambda^2a h'(a).
\end{equation}
Consequently,
\begin{equation}\label{eq:a-only-harmonic-family}
 \Delta_g(h\circ a)=0
 \quad\Longleftrightarrow\quad
 h(a)=C_0+C_1\arctan\!\left(\frac{\lambda a}{\mu}\right)
\end{equation}
for constants $C_0,C_1\in\C$.
\end{theorem}

\begin{proof}
The chain rule for the Laplace--Beltrami operator and
\eqref{eq:hyperbolic-assignment-eigen-equation-II} give
\[
 \Delta_g(h\circ a)
 =h''(a)|\nabla a|_g^2+h'(a)\Delta_ga,
\]
which is \eqref{eq:a-only-laplacian}.  If it vanishes and $p=h'$, then
\[
 (\mu^2+\lambda^2a^2)p'(a)+2\lambda^2a p(a)=0.
\]
Thus
\[
 p(a)=\frac{C}{\mu^2+\lambda^2a^2},
\]
and integration yields
\[
 h(a)=C_0+\frac{C}{\mu\lambda}
       \arctan\!\left(\frac{\lambda a}{\mu}\right).
\]
Absorbing the nonzero constant $(\mu\lambda)^{-1}$ into $C_1$ proves the
classification on $I$.  Here no values of $h$ are left untested: the assignment
$a=-x/y$ maps $\Omega_I$ onto $I$.  The converse follows by differentiation.
\end{proof}

\begin{corollary}[A half-line harmonic measure]
\label{cor:a-half-line-harmonic-measure}
The bounded function
\begin{equation}\label{eq:a-half-line-harmonic-measure}
 \omega_-(a)
 =\frac12+\frac1\pi
  \arctan\!\left(\frac{\lambda a}{\mu}\right)
\end{equation}
is the Poisson extension of the almost-everywhere boundary datum
$\mathbf 1_{(-\infty,0)}$ in the $x$-coordinate.  Equivalently,
\begin{equation}\label{eq:a-angle-identity}
 \omega_-(a)=\frac{\operatorname{Arg} w}{\pi},
 \qquad 0<\operatorname{Arg} w<\pi.
\end{equation}
\end{corollary}

\begin{proof}
Writing $w=re^{i\theta}$ with $0<\theta<\pi$, one has
\[
 \frac{\lambda a}{\mu}=-\frac{X}{y}=-\cot\theta.
\]
On this range,
$\arctan(-\cot\theta)=\theta-\pi/2$, which proves
\eqref{eq:a-angle-identity}.  The angle $\theta/\pi$ tends to zero on the
positive half-line and to one on the negative half-line.  It is bounded and
harmonic.  Directly, with $c=\mu/\lambda$,
\begin{align*}
 \int_{-\infty}^0P_y^{\mu,\lambda}(x-\xi)\,d\xi
 &=\frac1\pi\left[
   \arctan\!\left(\frac{\xi-x}{cy}\right)
   \right]_{\xi=-\infty}^{\xi=0}\\
 &=\frac12-\frac1\pi\arctan\!\left(\frac{x}{cy}\right)
  =\frac12+\frac1\pi\arctan\!\left(\frac{\lambda a}{\mu}\right).
\end{align*}
This identifies it as the stated harmonic measure without imposing continuity
at the jump point.
\end{proof}

The holomorphic function behind this corollary is the principal logarithm on
the upper half-plane:
\begin{equation}\label{eq:log-w-assignment}
 \operatorname{Log} w
 =\log y+\frac12\log\left(1+\frac{\lambda^2a^2}{\mu^2}\right)
 +i\left[
   \frac\pi2+\arctan\!\left(\frac{\lambda a}{\mu}\right)
 \right].
\end{equation}
The example lies outside the continuous compactified-boundary and finite-energy
classes.  Its boundary datum jumps at $x=0$ and has incompatible limits along
the two ends that meet at $\infty$.  Moreover, in polar coordinates for $w$,
\[
 \int_{\HH}|\nabla(\operatorname{Arg} w)|^2\,dX\,dy
 =\int_0^\pi\int_0^\infty\frac1{r^2}r\,dr\,d\theta
 =\infty.
\]
Thus the half-line harmonic measure is a valid bounded Poisson solution for
$L^\infty$ data, but not a member of the energy theorem above.  These domain
distinctions also explain why the unbounded assignment $a=-x/y$ itself cannot
serve as ideal-boundary Dirichlet data.
